% TOP_PROBLEM: {"id": 3040, "title": "Monotone 4-term Arithmetic Progressions", "category_id": 2, "category": {"id": 2, "name": "combinatorics", "display_name": "Combinatorics", "description": "Counting problems, graph theory, discrete structures.", "slug": "combinatorics", "order_index": 2, "created_at": "2026-07-31T15:26:25.671Z"}, "status": "open", "problem_number": "OPG-618", "set_id": 12}
% FIRST_POSTED: 2026-10-01
% ENTRY_KIND: statement_only
% UnsolvedMath ID 3040; source catalogue code OPG-618
% !TeX program = lualatex
% Definition and sources only; this file is not an LLM solution attempt.
\documentclass[11pt,a4paper]{article}
\usepackage[margin=25mm]{geometry}
\usepackage{fontspec}
\setmainfont{lmroman10-regular.otf}[BoldFont=lmroman10-bold.otf,
 ItalicFont=lmroman10-italic.otf,BoldItalicFont=lmroman10-bolditalic.otf]
\usepackage{amsmath,amssymb,mathtools}
\usepackage{unicode-math}
\setmathfont{latinmodern-math.otf}
\AtBeginDocument{\let\setminus\smallsetminus}
\usepackage{xurl,hyperref}
\hypersetup{colorlinks=true,urlcolor=blue}
\setcounter{secnumdepth}{0}
\setlength{\parindent}{0pt}
\setlength{\parskip}{5pt}
\setlength{\emergencystretch}{3em}
\begin{document}
\section*{Monotone 4-term Arithmetic Progressions}\label{3040}
{\small UnsolvedMath ID 3040 · OPG-618 · Combinatorics · OpenGarden}
\subsection{Definitions and mathematical statement}
Let $\mathbb N=\{1,2,\ldots\}$ and let $\pi:\mathbb N\to\mathbb N$ be a
bijection. For positive integers $a,d$, the four indices
$a,a+d,a+2d,a+3d$ form a nonconstant arithmetic progression. Call this
progression monotone for $\pi$ if either
\[
 \pi(a)<\pi(a+d)<\pi(a+2d)<\pi(a+3d)
\]
or all three inequalities are reversed. All inequalities are strict,
as the values of a permutation are distinct.

The question is whether every permutation $\pi$ admits such $a,d$.
Equivalently, regard a permutation as a one-sided sequence
$b_1,b_2,\ldots$ listing every positive integer exactly once. Must four
terms occurring in their sequence order form an increasing or decreasing
arithmetic progression? To relate these formulations, let $\pi(t)$ be the
position at which the integer $t$ occurs; the inverse of any permutation is
again a permutation.

The positions of the four terms in the sequence formulation need not be
equally spaced. In the displayed formulation the indices are equally
spaced and the values are merely monotone, not necessarily an arithmetic
progression. This distinction prevents imposing an extra requirement absent
from the problem. The ordering has order type $\mathbb N$, so arrangements
indexed by all integers do not automatically answer it. A counterexample
must be a single bijection avoiding every increasing and every decreasing
four-term arithmetic progression in this sense.
\subsection{Short English statement}
Does every one-sided ordering of the positive integers contain four arithmetic-progression values appearing in increasing or decreasing order?
\subsection{Sources}
\begin{itemize}
\item[S1] ULAM.AI and the UnsolvedMath Contributors, supplied problems.json, record 3040 (OPG-618), OpenGarden collection. Catalogue identity and source transcription; CC BY 4.0.
\url{https://huggingface.co/datasets/ulamai/UnsolvedMath}.
\item[S2] Open Problem Garden, Monotone 4-term Arithmetic Progressions. Original problem and discussion; the mathematical formulation below is expanded from the supplied source record.
\url{http://www.openproblemgarden.org/op/monotone_4_term_arithmetic_progressions}.
\end{itemize}
\end{document}
